Geocentric ephemerides of the Moon for September 2015 are given in the form
of a table. Each reading was taken at 00:00 UT.

\begin{center}
\begin{tabular}{c|ccc|ccc|c|c|c}
Date & \multicolumn{3}{c|}{R.A. ($\alpha$)}
  & \multicolumn{3}{c|}{Dec. ($\delta$)}
  & Angular Size ($\theta$) & Phase ($\phi$) & Elongation \\
 & h & m & s & $^\circ$ & $'$ & $''$ & $''$ & & Of Moon \\
\hline
Sep 01 & 0 & 36 & 46.02 & 3  & 6  & 16.8 & 1991.2 & 0.927 & $148.6^\circ$ W \\
Sep 02 & 1 & 33 & 51.34 & 7  & 32 & 26.1 & 1974.0 & 0.852 & $134.7^\circ$ W \\
Sep 03 & 2 & 30 & 45.03 & 11 & 25 & 31.1 & 1950.7 & 0.759 & $121.1^\circ$ W \\
Sep 04 & 3 & 27 & 28.48 & 14 & 32 & 4.3  & 1923.9 & 0.655 & $107.9^\circ$ W \\
Sep 05 & 4 & 23 & 52.28 & 16 & 43 & 18.2 & 1896.3 & 0.546 & $95.2^\circ$ W \\
\end{tabular}
\end{center}
% TABLE ABRIDGED — full data (Sep 01 to Sep 30) in crop edition; the complete
% table is also wired below as a figure.
\ioaafig{2016_DA_D2-f1.pdf}

The composite graphic\footnote{Credit: NASA's Scientific Visualization
Studio} below shows multiple snapshots of the Moon taken at different times
during the total lunar eclipse, which occurred in this month. For each shot,
the centre of frame was coinciding with the central north-south line of
umbra.

For this problem, assume that the observer is at the centre of the Earth and
angular size refers to angular diameter of the object / shadow.

\ioaafig{2016_DA_D2-f2.pdf}

\begin{description}
\item[(D2.1)] In September 2015, apogee of the lunar orbit is closest to

  New Moon / First Quarter / Full Moon / Third Quarter.

  Tick the correct answer in the Summary Answersheet. No justification for
  your answer is necessary. \hfill[3]
\item[(D2.2)] In September 2015, the ascending node of lunar orbit with
  respect to the ecliptic is closest to

  New Moon / First Quarter / Full Moon / Third Quarter.

  Tick the correct answer in the Summary Answersheet. No justification for
  your answer is necessary. \hfill[4]
\item[(D2.3)] Estimate the eccentricity, $e$, of the lunar orbit from the
  given data. \hfill[4]
\item[(D2.4)] Estimate the angular size of the umbra,
  $\theta_{\mathrm{umbra}}$, in terms of the angular size of the Moon,
  $\theta_{\mathrm{Moon}}$. Show your working on the image given on the
  backside of the Summary Answersheet. \hfill[8]
\item[(D2.5)] The angle subtended by the Sun at Earth on the day of the lunar
  eclipse is known to be $\theta_{\mathrm{Sun}} = 1915.0''$. In the figure
  below, $S_1 R_1$ and $S_2 R_2$ are rays coming from diametrically opposite
  ends of the solar disk. The figure is not to scale.

\ioaafig{2016_DA_D2-f3.pdf}

  Calculate the angular size of the penumbra, $\theta_{\mathrm{penumbra}}$,
  in terms of $\theta_{\mathrm{Moon}}$. Assume the observer to be at the
  centre of the Earth. \hfill[9]
\item[(D2.6)] Let $\theta_{\mathrm{Earth}}$ be angular size of the Earth as
  seen from the centre of the Moon. Calculate the angular size of the Moon,
  $\theta_{\mathrm{Moon}}$, as would be seen from the centre of the Earth on
  the eclipse day in terms of $\theta_{\mathrm{Earth}}$. \hfill[5]
\item[(D2.7)] Estimate the radius of the Moon, $R_{\mathrm{Moon}}$, in km
  from the results above. \hfill[3]
\item[(D2.8)] Estimate the shortest distance, $r_{\mathrm{perigee}}$, and the
  farthest distance, $r_{\mathrm{apogee}}$, to the Moon. \hfill[4]
\item[(D2.9)] Use appropriate data from September 10 to estimate the
  distance, $d_{\mathrm{Sun}}$, to the Sun from the Earth. \hfill[10]
\end{description}
